%% ----------------------------------------------------------------------------
%%  §3.5  Teoremi di confronto e altre proprietà dei limiti
%% ----------------------------------------------------------------------------
\section{Teoremi di confronto}
\label{sec:confronto}

Studiamo in questo paragrafo alcune relazioni tra limiti e ordinamento: che cosa si può
dedurre sul limite conoscendo il segno dei termini, e viceversa; e come si può
determinare il limite di una successione «intrappolandola» tra due successioni di cui
conosciamo il comportamento.

\subsection{Permanenza del segno}

\begin{teorema}[Permanenza del segno]\label{teo:permanenza-segno}\index{teorema!della permanenza del segno}
Se $\displaystyle\lim_{n\to+\infty}a_{n}=a>0$, esiste un numero $\nu$ tale che $a_{n}>0$
per ogni $n>\nu$.
\end{teorema}

\begin{schema}
\Ipotesi{$a_{n}\to a$, \ $a>0$.}
\Tesi{$a_{n}>0$ definitivamente.}
\end{schema}

\begin{strategia}
Se il limite è positivo, i termini della successione finiscono per stare in un intorno
di $a$; basta scegliere l'intorno abbastanza piccolo da non contenere lo zero. Il raggio
$\epsilon=a/2$ va bene: l'intorno $(a/2,\,3a/2)$ sta tutto a destra dello zero. Anzi,
otteniamo un'informazione più precisa: definitivamente $a_{n}>a/2$, cioè i termini non
solo sono positivi, ma si mantengono lontani da zero.
\end{strategia}

\begin{dimostrazione}
Dato che $a>0$, possiamo scegliere $\epsilon=a/2$. Esiste quindi un numero $\nu$ per cui
$\abs{a_{n}-a}<a/2$ per ogni $n>\nu$. Ciò equivale a $-a/2<a_{n}-a<a/2$. In particolare
abbiamo
\begin{equation}\label{eq:permanenza-passo}
  a_{n}>a-\frac a2=\frac a2>0 \qquad\perogni n>\nu. \qedhere
\end{equation}
\end{dimostrazione}

Sottolineiamo che, se una successione $a_{n}$ converge a un numero reale positivo $a$,
\emph{non} si può affermare in generale che \emph{tutti} i termini della successione
siano positivi. Ad esempio, in analogia con la \eqref{eq:succ-n1n}, la successione
$a_{n}=(n-7)/n$ converge al numero $1$, però i primi termini della successione (da $a_{1}$
fino ad $a_{6}$) sono negativi e $a_{7}=0$; il numero $\nu$ nella tesi del teorema della
permanenza del segno, in questo caso, può essere scelto uguale a $7$, e $7$ è il più
piccolo valore possibile.

\begin{osservazione}
Con la stessa dimostrazione (scegliendo $\epsilon=a-L$) si prova la versione più
generale: \emph{se $a_{n}\to a$ e $a>L$, allora definitivamente $a_{n}>L$}; e
analogamente, se $a<L$, definitivamente $a_{n}<L$. Il caso $L=0$ è il teorema.
\end{osservazione}

\begin{corollario}\label{cor:limite-non-negativo}
Se $\displaystyle\lim_{n\to+\infty}a_{n}=a$, e se $a_{n}\ge0$ per ogni $n$, allora anche
$a\ge0$.
\end{corollario}

\begin{schema}
\Ipotesi{$a_{n}\to a\in\R$; \ $a_{n}\ge0$ per ogni $n$.}
\Tesi{$a\ge0$.}
\end{schema}

\begin{strategia}
Ragioniamo per assurdo: se il limite fosse negativo, la permanenza del segno (applicata
a $-a_{n}$) renderebbe i termini definitivamente negativi, contro l'ipotesi.
\end{strategia}

\begin{dimostrazione}
Se per assurdo fosse $a<0$, il teorema della permanenza del segno, applicato alla
successione $-a_{n}$ (che converge a $-a>0$ per il \cref{teo:operazioni-limiti}),
comporterebbe che $-a_{n}>0$, cioè $a_{n}<0$, per $n$ grande, contro l'ipotesi.
\end{dimostrazione}

\begin{corollario}[Confronto dei limiti]\label{cor:confronto-limiti}
Se $\displaystyle\lim_{n\to+\infty}a_{n}=a$, $\displaystyle\lim_{n\to+\infty}b_{n}=b$, e
se $a_{n}\ge b_{n}$ per ogni $n$, allora $a\ge b$.
\end{corollario}

\begin{schema}
\Ipotesi{$a_{n}\to a\in\R$, \ $b_{n}\to b\in\R$; \ $a_{n}\ge b_{n}$ per ogni $n$.}
\Tesi{$a\ge b$.}
\end{schema}

\begin{strategia}
Confrontare due successioni equivale a confrontare la loro differenza con zero: basta
applicare il corollario precedente ad $a_{n}-b_{n}$, il cui limite è noto per il
\cref{teo:operazioni-limiti}.
\end{strategia}

\begin{dimostrazione}
Per ottenere la dimostrazione di quest'ultimo corollario basta applicare il corollario
precedente alla successione $a_{n}-b_{n}\ge0$, che converge ad $a-b$.
\end{dimostrazione}

Possiamo schematizzare i risultati ottenuti nel modo seguente (si noti la differenza tra
i segni $>$ e $\ge$):
\begin{align}
  a_{n}\to a,\ \ a>0\ \ &\implica\ \ a_{n}>0\ \text{ definitivamente}; \label{eq:schema-perm1}\\
  a_{n}\to a,\ \ a_{n}\ge0\ \ \perogni n\in\N\ \ &\implica\ \ a\ge0; \label{eq:schema-perm2}\\
  a_{n}\to a,\ \ b_{n}\to b,\ \ a_{n}\ge b_{n}\ \ \perogni n\in\N\ \ &\implica\ \ a\ge b. \label{eq:schema-perm3}
\end{align}

\begin{attenzione}[Le disuguaglianze strette non passano al limite]
Nelle \eqref{eq:schema-perm2} e \eqref{eq:schema-perm3} non si può sostituire $\ge$ con
$>$ nella tesi, nemmeno se l'ipotesi è stretta: $a_{n}=1/n>0$ per ogni $n$, ma il limite
è $0$, che non è positivo. Passando al limite, «$>$» si indebolisce in «$\ge$». Al
contrario, nella \eqref{eq:schema-perm1} l'ipotesi deve essere stretta: se fosse $a=0$, i
termini potrebbero avere qualsiasi segno, come accade per $(-1)^{n}/n$.
\end{attenzione}

\subsection{Il teorema dei carabinieri}

\begin{teorema}[dei carabinieri]\label{teo:carabinieri}\index{teorema!dei carabinieri}
Siano $a_{n}$, $b_{n}$, $c_{n}$ tre successioni tali che
\begin{equation}\label{eq:carabinieri-ipotesi}
  a_{n}\le c_{n}\le b_{n} \qquad\perogni n\in\N .
\end{equation}
Se $\displaystyle\lim_{n\to+\infty}a_{n}=\lim_{n\to+\infty}b_{n}=a$, allora anche la
successione $c_{n}$ è convergente e $\displaystyle\lim_{n\to+\infty}c_{n}=a$.
\end{teorema}

\begin{schema}
\Ipotesi{$a_{n}\le c_{n}\le b_{n}$ per ogni $n$ (basta definitivamente); \
$a_{n}\to a$, \ $b_{n}\to a$ (lo \emph{stesso} limite).}
\Tesi{$c_{n}\to a$ \ (in particolare, $c_{n}$ ammette limite).}
\end{schema}

\begin{strategia}
Si noti che la tesi contiene un'affermazione di \emph{esistenza}: non si suppone che
$c_{n}$ abbia limite, lo si dimostra. Non possiamo quindi usare il
\cref{cor:confronto-limiti}, che presuppone l'esistenza dei limiti; dobbiamo tornare alla
definizione. L'idea è visiva (\cref{fig:carabinieri}): fissato un intorno
$(a-\epsilon,a+\epsilon)$, definitivamente vi cadono sia $a_{n}$ sia $b_{n}$; e poiché
$c_{n}$ è «stretto» tra i due, ci cade anche lui. In formule, basta leggere le due
disuguaglianze $\abs{a_{n}-a}<\epsilon$ e $\abs{b_{n}-a}<\epsilon$ nella forma «doppia
disuguaglianza» e concatenarle con la \eqref{eq:carabinieri-ipotesi}.
\end{strategia}

\begin{dimostrazione}
Per ipotesi, per ogni $\epsilon>0$,
\begin{equation}\label{eq:carabinieri-passo1}
  \esiste\nu_{1}:\ \abs{a_{n}-a}<\epsilon\quad\perogni n>\nu_{1}; \qquad
  \esiste\nu_{2}:\ \abs{b_{n}-a}<\epsilon\quad\perogni n>\nu_{2}.
\end{equation}
Ricordiamo che le disuguaglianze con il valore assoluto si possono anche scrivere
\begin{equation}\label{eq:carabinieri-passo2}
  a-\epsilon<a_{n}<a+\epsilon; \qquad a-\epsilon<b_{n}<a+\epsilon .
\end{equation}
Quindi, se $n>\nu=\max\{\nu_{1},\nu_{2}\}$, risulta
\begin{equation}\label{eq:carabinieri-passo3}
  a-\epsilon<a_{n}\le c_{n}\le b_{n}<a+\epsilon .
\end{equation}
Perciò $\abs{c_{n}-a}<\epsilon$ per ogni $n>\nu$, come volevasi dimostrare.
\end{dimostrazione}

\begin{figure}[htb]
\centering
\begin{tikzpicture}
\begin{axis}[width=0.92\textwidth, height=5.2cm, xmin=0, xmax=31, ymin=-0.1, ymax=2.1,
  xtick={5,10,15,20,25,30}, ytick={1}, yticklabels={$a$}, xlabel={$n$}]
  \fill[verde!15] (axis cs:0,0.85) rectangle (axis cs:31,1.15);
  \addplot[tratteggio, domain=0:31] {1};
  \addplot[only marks, mark=triangle*, mark size=1.8pt, rame, samples at={1,...,30}] {1-1/x};
  \addplot[only marks, mark=triangle*, mark size=1.8pt, rame, mark options={rotate=180}, samples at={1,...,30}] {1+1/x};
  \addplot[only marks, mark=*, mark size=1.3pt, blu, samples at={1,...,30}] {1+0.8*sin(deg(1.3*x))/x};
  \node[font=\scriptsize, rame, anchor=west] at (axis cs:1.6,0.22) {$a_{n}$};
  \node[font=\scriptsize, rame, anchor=west] at (axis cs:1.6,1.93) {$b_{n}$};
  \node[font=\scriptsize, blu, anchor=west] at (axis cs:1.6,1.72) {$c_{n}$};
  \node[font=\scriptsize, verde!50!black, anchor=east] at (axis cs:30.5,1.25) {$(a-\epsilon,a+\epsilon)$};
\end{axis}
\end{tikzpicture}
\caption{Il teorema dei carabinieri: la successione $c_{n}$ (in blu) è compresa tra
$a_{n}$ e $b_{n}$ (i «carabinieri», in rame), che convergono allo stesso limite $a$.
Quando i due carabinieri sono entrati nella striscia verde $(a-\epsilon,a+\epsilon)$,
anche $c_{n}$, che sta in mezzo, è costretta a entrarvi.}
\label{fig:carabinieri}
\end{figure}

\begin{osservazione}
Nell'enunciato basta che la \eqref{eq:carabinieri-ipotesi} valga definitivamente, cioè
per $n$ maggiore di un certo $\nu_{3}$: nella dimostrazione si prende allora
$\nu=\max\{\nu_{1},\nu_{2},\nu_{3}\}$. In generale, il comportamento di una successione
per $n$ grande non dipende dai primi termini: modificando un numero finito di termini il
limite (se esiste) non cambia. Il nome del teorema, d'uso comune in Italia, evoca un
arrestato condotto tra due carabinieri: dove vanno loro, deve andare anche lui. In altri
paesi si chiama «teorema del sandwich» o «del confronto».
\end{osservazione}

\begin{idea}[Le inviluppanti di un'oscillazione smorzata]
Un esempio fisico: lo spostamento di una struttura che vibra con ampiezza decrescente,
campionato a istanti successivi, è una successione $c_{n}$ che oscilla attorno alla
posizione di equilibrio; ma se l'ampiezza $A_{n}$ tende a zero, le due curve
\emph{inviluppanti} $-A_{n}$ e $A_{n}$ fanno da carabinieri, e lo spostamento tende a
zero qualunque sia la forma dettagliata delle oscillazioni. È esattamente la situazione
della \cref{fig:carabinieri}, e del \cref{teo:limitata-infinitesima} qui sotto.
\end{idea}

Valgono analoghi teoremi di confronto anche per i limiti infiniti; qui basta un solo
«carabiniere».

\begin{teorema}[Confronto per limiti infiniti]\label{teo:confronto-infinito}
Se $a_{n}\le b_{n}$ per ogni $n\in\N$, allora
\begin{align}
  a_{n}\to+\infty &\ \implica\ b_{n}\to+\infty; \label{eq:confronto-piu-inf}\\
  b_{n}\to-\infty &\ \implica\ a_{n}\to-\infty. \label{eq:confronto-meno-inf}
\end{align}
\end{teorema}

\begin{schema}
\Ipotesi{$a_{n}\le b_{n}$ per ogni $n$ (basta definitivamente); \ $a_{n}\to+\infty$
(risp.\ $b_{n}\to-\infty$).}
\Tesi{$b_{n}\to+\infty$ (risp.\ $a_{n}\to-\infty$).}
\end{schema}

\begin{strategia}
Basta la definizione: se la successione più piccola supera definitivamente un livello
$M$, la più grande lo supera a maggior ragione, con lo stesso indice $\nu$. Non occorre
un secondo carabiniere, perché verso $+\infty$ non c'è nulla da «stringere».
\end{strategia}

\begin{dimostrazione}
Dimostriamo la \eqref{eq:confronto-piu-inf} (la prova della \eqref{eq:confronto-meno-inf}
è analoga). Per ipotesi $a_{n}\to+\infty$, che per la definizione
\eqref{eq:def-piu-inf} significa
\begin{equation}\label{eq:confronto-inf-passo}
  \perogni M>0\ \ \esiste\nu :\ a_{n}>M \quad\perogni n>\nu .
\end{equation}
Dato che $b_{n}\ge a_{n}$ per ogni $n\in\N$, si ottiene la tesi:
$b_{n}\ge a_{n}>M$ per ogni $n>\nu$.
\end{dimostrazione}

\subsection{Successioni infinitesime}

Riportiamo infine due ulteriori proprietà dei limiti di successioni; in particolare la
seconda, enunciata sotto forma di teorema, è importante per le applicazioni.

\begin{proposizione}\label{prop:infinitesima-modulo}
La successione $a_{n}$ converge a zero se e soltanto se $\abs{a_{n}}$ converge a zero.
\end{proposizione}

\begin{schema}
\Ipotesi{$a_{n}$ successione di numeri reali.}
\Tesi{$a_{n}\to0$ \ $\sse$ \ $\abs{a_{n}}\to0$.}
\end{schema}

\begin{strategia}
Quando il limite è zero, la definizione coinvolge soltanto la distanza da zero, cioè il
valore assoluto: scritte per esteso, le due condizioni sono la stessa frase, perché
$\bigl|\abs{a_{n}}\bigr|=\abs{a_{n}}$.
\end{strategia}

\begin{dimostrazione}
Posto $b_{n}=\abs{a_{n}}$, in base alla definizione \eqref{eq:def-limite} $b_{n}$
converge a zero se e solo se
\begin{equation}\label{eq:modulo-zero-passo}
  \perogni\epsilon>0\ \ \esiste\nu :\ \abs{b_{n}}<\epsilon \quad\perogni n>\nu .
\end{equation}
Dato che $\abs{b_{n}}=\bigl|\abs{a_{n}}\bigr|=\abs{a_{n}}$ per ogni $n\in\N$, la
\eqref{eq:modulo-zero-passo} è equivalente alla convergenza a zero della successione
$a_{n}$.
\end{dimostrazione}

Si noti che, nella proposizione precedente, è importante considerare non solo
successioni convergenti, ma più in particolare successioni \emph{convergenti a zero}. Ad
esempio, se $a_{n}=(-1)^{n}$, allora $a_{n}$ non è convergente, mentre
$b_{n}=\abs{a_{n}}$ è la successione costante $b_{n}=1$, che ovviamente converge a $1$. In
generale vale solo l'implicazione del \cref{lem:lim-modulo}: se $a_{n}\to a$ allora
$\abs{a_{n}}\to\abs a$.

\begin{teorema}[Prodotto di una limitata per un'infinitesima]\label{teo:limitata-infinitesima}\index{successione!infinitesima}
Se $a_{n}$ è una successione limitata e $b_{n}$ è una successione che converge a zero,
allora la successione prodotto $a_{n}\cdot b_{n}$ converge a zero.
\end{teorema}

\begin{schema}
\Ipotesi{$\abs{a_{n}}\le M$ per ogni $n$; \ $b_{n}\to0$.}
\Tesi{$a_{n}b_{n}\to0$.}
\end{schema}

\begin{strategia}
Si noti che $a_{n}$ \emph{non} è supposta convergente (potrebbe essere $(-1)^{n}$): il
\cref{teo:operazioni-limiti} non si applica. Ma $\abs{a_{n}b_{n}}\le M\abs{b_{n}}$, e la
successione a destra tende a zero: possiamo quindi usare i carabinieri $\pm M\abs{b_{n}}$
(primo metodo), oppure verificare direttamente la definizione nella forma
\eqref{eq:def-limite-c} (secondo metodo).
\end{strategia}

\begin{dimostrazione}[Primo metodo]
Per l'ipotesi di limitatezza si ha
\begin{equation}\label{eq:lim-inf-passo1}
  \abs{a_{n}b_{n}}=\abs{a_{n}}\cdot\abs{b_{n}}\le M\abs{b_{n}} \qquad\perogni n\in\N,
\end{equation}
che, per la proprietà \eqref{eq:mod-le} del valore assoluto, equivale a
\begin{equation}\label{eq:lim-inf-passo2}
  -M\abs{b_{n}}\le a_{n}b_{n}\le M\abs{b_{n}} \qquad\perogni n\in\N .
\end{equation}
Dato che per ipotesi $b_{n}\to0$, per la proposizione precedente anche la successione
$\abs{b_{n}}$ converge a zero, e quindi $\pm M\abs{b_{n}}\to0$. Per il teorema dei
carabinieri, dalla \eqref{eq:lim-inf-passo2} si deduce infine che $a_{n}b_{n}\to0$.
\end{dimostrazione}

\begin{dimostrazione}[Secondo metodo]
Per la definizione di limite si ha
$\perogni\epsilon>0\ \esiste\nu:\ \abs{b_{n}}<\epsilon$ per ogni $n>\nu$. Dall'ipotesi di
limitatezza si ottiene poi
\begin{equation}\label{eq:lim-inf-passo3}
  \abs{a_{n}b_{n}}=\abs{a_{n}}\cdot\abs{b_{n}}\le M\abs{b_{n}}<M\epsilon \qquad\perogni n>\nu,
\end{equation}
che equivale (si veda la \eqref{eq:def-limite-c}) al fatto che la successione prodotto
$a_{n}b_{n}$ converge a zero.
\end{dimostrazione}

\begin{esempio}
A titolo di esempio verifichiamo che
\begin{equation}\label{eq:es-limitata-inf}
  \lim_{n\to+\infty}\frac{(-1)^{n}(n+5)}{3n^{2}+1}=0;
\end{equation}
infatti si tratta del limite del prodotto della successione limitata $a_{n}=(-1)^{n}$
per la successione infinitesima
\begin{equation}\label{eq:es-limitata-inf-b}
  b_{n}=\frac{n+5}{3n^{2}+1}=\frac{1/n+5/n^{2}}{3+1/n^{2}}\to\frac03=0 .
\end{equation}
\end{esempio}

\begin{esempio}\label{es:senn-n}
Come ulteriore esempio verifichiamo che
\begin{equation}\label{eq:lim-senn-n}
  \lim_{n\to+\infty}\frac{\sen n}{n}=0 .
\end{equation}
Il limite è zero perché è il limite del prodotto della successione limitata
$a_{n}=\sen n$ ($\abs{a_{n}}=\abs{\sen n}\le1$ per ogni $n\in\N$, per la
\eqref{eq:sen-cos-limitate}) per la successione infinitesima $b_{n}=1/n$. Il lettore non
confonda questo limite con quello che dimostreremo nel \cref{teo:sen-x-su-x}, uguale a
$1$: là la successione $a_{n}$ converge a zero, mentre la successione in esame nella
\eqref{eq:lim-senn-n} è $\sen a_{n}/a_{n}$ con $a_{n}=n\to+\infty$.
\end{esempio}
